\documentclass[10pt,a4paper]{article}
\usepackage[utf8]{inputenc}
\usepackage[T1]{fontenc}
\usepackage{lmodern}
\usepackage[spanish,es-noshorthands]{babel}
\usepackage[margin=1.9cm,top=2.2cm,bottom=2cm]{geometry}
\usepackage{booktabs,tabularx,array,enumitem,fancyhdr,titlesec,xcolor}
\usepackage{pgfplots}
\pgfplotsset{compat=1.16}
\usepackage[hidelinks]{hyperref}
\hypersetup{pdftitle={Reporte de Sprints y métricas ágiles - Corte 2},
 pdfauthor={DataForge Consulting},pdfsubject={Cierre del caso S1/S2 e incremento integrado; 16 de octubre de 2026}}
\newcommand{\repourl}{https://github.com/OscarSantiagoMerino/DataForge-Consulting}
\definecolor{forge}{HTML}{0B4F71}
\definecolor{teal}{HTML}{12908E}
\definecolor{alert}{HTML}{B5123B}
\pagestyle{fancy}\fancyhf{}
\lhead{\small Reporte de Sprints, Corte 2. FinBank S.A.}
\rhead{\small DataForge Consulting}
\cfoot{\small\thepage}
\setlength{\headheight}{13pt}
\titleformat{\section}{\normalsize\bfseries\color{forge}}{\thesection.}{0.5em}{}
\titlespacing*{\section}{0pt}{9pt}{4pt}
\setlength{\parindent}{0pt}\setlength{\parskip}{4pt}
\setlist{nosep,leftmargin=1.2em}
\newcolumntype{L}{>{\raggedright\arraybackslash}X}
\renewcommand{\arraystretch}{1.12}

\begin{document}
\begin{center}
{\Large\bfseries Reporte de Sprints y métricas ágiles}\\[2pt]
Programa de Analítica Predictiva: Sprints 1 y 2 con Scrum y XP\\[2pt]
{\small Cliente: FinBank S.A., Vicepresidencia de Riesgos y Analítica \quad Código: GMAC-FB-2026-01}\\
{\small Daniel Felipe Martínez Hernández \,·\, Rafael Alejandro Baracaldo Doikov \,·\, Oscar Santiago Merino Suárez}
\end{center}
\vspace{-6pt}\rule{\linewidth}{0.4pt}

\section{Alcance del reporte}
En C1 gestionamos el caso con PMBOK y un tablero Kanban. En C2 retomamos los mismos dos sprints,
el backlog y el bloqueo de acceso, dentro de las 12 semanas y COP \$180.000.000 del Charter.
Scrum aporta objetivos, eventos y medición de capacidad; XP incorpora parejas, pruebas y revisión del código.
Rafael Baracaldo actúa como Product Owner, Oscar Merino como Scrum Master y los tres como Developers.
Daniel Martínez conserva el reporte formal al patrocinador.
Las fechas y actas de gestión describen el \textbf{escenario académico}; los commits, PR y ejecuciones
de octubre sustentan la evidencia técnica. Actualización: 16 de octubre de 2026.

\section{Cierre de los sprints}
{\small
\begin{tabularx}{\linewidth}{@{}l l L c c l@{}}
\toprule
Sprint & Fechas del caso & Objetivo & Compromiso & Cierre & Resultado \\
\midrule
S1 & 10--21 ago. & Establecer gobierno, acceso y perfilamiento. & 16 & 16 & Cumplido \\
S2 & 24 ago.--4 sep. & Preparar la base de datos de mora y churn. & 20 & 15 & Parcial \\
\bottomrule
\end{tabularx}

\vspace{5pt}
\begin{tabularx}{\linewidth}{@{}l L c l@{\hspace{14pt}}l L c l@{}}
\toprule
\multicolumn{4}{@{}l}{\textbf{Backlog del Sprint 1}} & \multicolumn{4}{l}{\textbf{Backlog del Sprint 2}}\\
ID & Ítem & Pts & Estado & ID & Ítem & Pts & Estado \\
\midrule
US-01 & Charter aprobado en el caso & 3 & Hecho & T-05.0 & Muestra sintética y layout de mora & 2 & Hecho \\
US-02 & Acuerdo y condiciones de acceso & 5 & Hecho & T-05.6 & Reglas de validación y pruebas & 3 & Hecho \\
US-03 & Perfilamiento de la muestra & 5 & Hecho & T-05.1 & Reporte de calidad & 4 & Hecho \\
US-04 & Tablero, repositorio y DoD & 3 & Hecho & T-05.4 & Etiquetas de mora a 30 y 60 días & 3 & Hecho \\
 & & & & T-05.2 & Diccionario v1: 18 variables & 3 & Hecho \\
 & & & & T-05.5 & Extracto histórico de churn & 5 & \textcolor{alert}{Bloqueado} \\
\bottomrule
\end{tabularx}}

El incremento técnico está integrado en \texttt{main} por el PR \#8: genera 5.000 registros sintéticos,
valida la entrada y calcula completitud, diccionario y etiquetas. Excluye cuatro de 22 variables con
completitud inferior al 80\%; conserva 18 y pasa 16 pruebas. Las tasas de mora son 5,85\% a 30 días y
2,37\% a 60 sobre 4.681 observados; 319 etiquetas conservan NaN. M2 mantiene pendiente churn.
El registro de aceptación del incremento se completará durante la Review.

\section{Velocidad y burndown}
\begin{center}
\begin{tikzpicture}
\begin{axis}[width=5.3cm,height=4.3cm,ybar,bar width=7pt,ymin=0,ymax=24,
 symbolic x coords={S1,S2},xtick=data,enlarge x limits=0.5,
 ylabel={Puntos de historia},ylabel style={font=\scriptsize},tick label style={font=\scriptsize},
 nodes near coords,nodes near coords style={font=\tiny},
 legend style={font=\tiny,at={(0.5,1.03)},anchor=south,legend columns=2,draw=none},axis lines*=left]
\addplot[fill=gray!45,draw=gray] coordinates {(S1,16) (S2,20)};
\addplot[fill=forge,draw=forge] coordinates {(S1,16) (S2,15)};
\legend{Compromiso,Cierre}
\end{axis}
\end{tikzpicture}\hfill
\begin{tikzpicture}
\begin{axis}[width=6.1cm,height=4.3cm,xmin=0,xmax=10,ymin=0,ymax=21,xtick={0,2,...,10},
 xlabel={Sprint 1, día hábil},label style={font=\scriptsize},tick label style={font=\scriptsize},
 legend style={font=\tiny,draw=none,at={(0.5,1.03)},anchor=south,legend columns=2},axis lines*=left]
\addplot[dashed,gray,thick] coordinates {(0,16) (10,0)};
\addplot[const plot,forge,thick,mark=*,mark size=1pt] coordinates
 {(0,16) (1,16) (2,16) (3,13) (4,13) (5,8) (6,8) (7,5) (8,5) (9,5) (10,0)};
\legend{Ideal,Reconstruido}
\end{axis}
\end{tikzpicture}\hfill
\begin{tikzpicture}
\begin{axis}[width=6.1cm,height=4.3cm,xmin=0,xmax=10,ymin=0,ymax=21,xtick={0,2,...,10},
 xlabel={Sprint 2, día hábil},label style={font=\scriptsize},tick label style={font=\scriptsize},
 legend style={font=\tiny,draw=none,at={(0.5,1.03)},anchor=south,legend columns=2},axis lines*=left]
\addplot[dashed,gray,thick] coordinates {(0,20) (10,0)};
\addplot[const plot,forge,thick,mark=*,mark size=1pt] coordinates
 {(0,20) (1,20) (2,20) (3,18) (4,18) (5,15) (6,15) (7,15) (8,8) (9,8) (10,5)};
\addplot[alert,thick,dotted] coordinates {(6,5) (10,5)};
\node[font=\tiny,text=alert,anchor=north east] at (axis cs:10,4.6) {5 pts bloqueados};
\legend{Ideal,Reconstruido}
\end{axis}
\end{tikzpicture}
\end{center}
{\footnotesize Curvas reconstruidas con los cierres diarios del caso. Fuente: \texttt{docs/SPRINT\_METRICS.md}
y \texttt{docs/evidence/sprint\_burndown.csv}. R-01 y T-05.3 están incluidos en las historias, sin puntos adicionales.}

\newpage
{\small
\begin{tabularx}{\linewidth}{@{}l c c c L@{}}
\toprule
Sprint & Compromiso & Cierre & Cumplimiento & Lectura \\
\midrule
S1 & 16 & 16 & 100\% & US-04 cierra en D3; US-03 en D5; US-01 en D7 y US-02 en D10. \\
S2 & 20 & 15 & 75\% & Descuenta 2 puntos en D3, 3 en D5, 7 en D8 y 3 en D10; churn conserva 5. \\
\midrule
Velocidad media & & 15,5 & & Referencia inicial para el siguiente Planning. \\
\bottomrule
\end{tabularx}}

El compromiso creció 25\% entre S1 y S2. Los 15 puntos cerrados dejan una referencia de capacidad de
15 a 16 por sprint. T-05.5 depende de una credencial externa; PO y Developers deberán comprobar esa
dependencia antes de comprometerla. Las curvas concilian la bitácora con los puntos de cada ítem.

\section{Bitácora de Daily Scrum}
Daily previsto de quince minutos. La selección recoge los días en que cambió el plan del caso.
Los registros originales de reuniones se acreditarán según el alcance del simulacro admitido por el docente.

{\footnotesize
\begin{tabularx}{\linewidth}{@{}l L L L@{}}
\toprule
Día & Hecho desde el último Daily & Plan del día & Impedimento o decisión \\
\midrule
S1 D1 & Planning: cuatro ítems, 16 puntos. & Daniel prepara Charter; Oscar el tablero; Rafael pide la muestra. & Solicitar acceso y muestra de apoyo. \\
S1 D3 & Tablero y DoD disponibles (US-04); muestra recibida. & Perfilar la base y tramitar el acuerdo. & Acceso a bodega pendiente (R1). \\
S1 D5 & Perfilamiento cerrado (US-03). & Revisar Charter con patrocinador. & Centralizar el script de perfilamiento. \\
S1 D7 & Charter aprobado en el caso (US-01). & Atender observaciones de protección de datos. & Acuerdo pendiente de firma en el caso. \\
S1 D10 & Acuerdo cerrado en el caso (US-02); 16 puntos. & Review, retro y Planning S2. & Adoptar repositorio y CI. \\
\midrule
S2 D1 & Planning: seis ítems, 20 puntos. & Preparar mora y solicitar acceso a churn. & Dependencia externa abierta. \\
S2 D3 & Muestra y layout listos (T-05.0). & P1: Rafael conduce y Daniel revisa reglas y pruebas. & Revisar faltantes y columnas personales. \\
S2 D5 & Reglas cerradas (T-05.6); siete pruebas de la etapa pasan. & Reporte de calidad y refactor R-01. & Límites duplicados entre reglas y diccionario. \\
S2 D6 & Calidad en revisión; refactor documentado. & Intentar extracción de churn. & \textcolor{alert}{Credencial pendiente; T-05.5 bloqueado.} \\
S2 D7 & Scrum Master escala; PO reordena el backlog. & P2: Oscar conduce y Rafael revisa targets 30/60. & Credencial pendiente, día 2. \\
S2 D8 & Calidad aceptada y targets integrados en el caso. & Revisar las 18 variables del diccionario. & Churn sigue bloqueado, día 3. \\
S2 D10 & Diccionario cerrado en el caso; 15 puntos. & Review y retrospectiva. & T-05.5 vuelve al Product Backlog. \\
\bottomrule
\end{tabularx}}

\section{Actas de retrospectiva}
{\footnotesize
\begin{tabularx}{\linewidth}{@{}l L L L@{}}
\toprule
 & Empezar (Start) & Dejar de hacer (Stop) & Continuar (Continue) \\
\midrule
S1 & Centralizar scripts en el repo y ejecutar CI. Probar reglas antes del análisis y trabajar en parejas. & Mantener trabajo aislado en equipos personales. Aceptar resultados sin ejecución compartida. & Pedir una muestra de apoyo mientras se tramita acceso; permitió avanzar con R1 abierto. \\
S2 & Planear según capacidad. Verificar dependencias con una Definition of Ready y escalar bloqueos el mismo día. & Comprometer tareas atadas al calendario de otra área. Duplicar límites entre código y diccionario. & Rotar conductor y navegador. Mantener pruebas y revisión de cambios. Reordenar con el PO. \\
\bottomrule
\end{tabularx}}

\newpage
{\footnotesize
\begin{tabularx}{\linewidth}{@{}l L l l@{}}
\toprule
Origen & Acción de mejora & Responsable & Estado \\
\midrule
S1 & Repositorio, CI y DoD con pruebas y revisión por PR. & Oscar & Operativos \\
S1 & Registrar dos sesiones de parejas con rotación. & Rafael & P1/P2 en XP\_LOG \\
S2 & Usar 15--16 puntos como referencia de capacidad. & Rafael & Propuesta S3 \\
S2 & Verificar dependencias antes del compromiso (DoR). & Oscar & Para Planning S3 \\
S2 & Escalar impedimentos en 24 horas. & Oscar & Seguimiento en S3 \\
\bottomrule
\end{tabularx}}

\section{Prácticas XP}
{\footnotesize
\begin{tabularx}{\linewidth}{@{}l L L@{}}
\toprule
Práctica & Aplicación & Evidencia en el repositorio \\
\midrule
Programación en parejas & P1: Rafael/Daniel en reglas; P2: Oscar/Rafael en targets. Cambio de conductor documentado. & Coautoría en \texttt{211d6c3} y \texttt{b3bb92e}; \texttt{docs/XP\_LOG.md}, práctica 01. \\
Pruebas y TDD & La ausencia de \texttt{loan\_amount} expone un fallo. El refactor detecta esa columna. & \texttt{c90a325}: 1 falla, 11 pasan; \texttt{d9d2b52}: 12 pasan. Suite completa: 16; registros en \texttt{docs/evidence/tdd/}. \\
CI y propiedad colectiva & PR \#1 a \#8 aprobados e integrados; \#8 consolida los componentes en \texttt{main}. & CI exitoso en \texttt{0b53f59}; rama protegida. XP\_LOG distingue aportes y aprobaciones. \\
Diseño simple y refactor & Validación y diccionario consultan \texttt{schema.py}. Pasa de 99 líneas y cinco variables a 84 y 22. & Antes en \texttt{docs/refactoring/}; después en \texttt{src/finbank/validation.py}; pruebas del refactor conservadas. \\
\bottomrule
\end{tabularx}}

Repositorio: \url{\repourl}

\section{Diferencias de trabajo entre C1 y C2}
El Charter mantiene alcance, presupuesto y gobierno. Scrum organiza el compromiso por sprint y XP aporta
controles técnicos que se ejecutan junto con el código.

{\footnotesize
\begin{tabularx}{\linewidth}{@{}>{\raggedright\arraybackslash}p{2.5cm} L L@{}}
\toprule
 & C1: PMBOK y Kanban & C2: Scrum y XP \\
\midrule
Unidad de trabajo & Tarjeta que entra según el límite de trabajo en curso (WIP). & Sprint con objetivo, backlog priorizado y compromiso de puntos. \\
Roles & Líder, analista y facilitador; matriz RACI y responsable por tarjeta. & Product Owner, Scrum Master y Developers; colaboración sobre el código. \\
Métricas & Throughput, tiempo de ciclo, WIP y variaciones de costo y plazo. & Velocidad estimada frente al cierre y burndown diario. \\
Uso de las métricas & Ubicar esperas y controlar la línea base del proyecto. & Estimar capacidad y ajustar el compromiso del próximo sprint. \\
Impedimentos & Tarjeta bloqueada y escalamiento mediante control de cambios. & Daily detecta el bloqueo; SM escala y PO reordena el backlog. \\
Cierre del trabajo & Resultado reproducible, documentado y validado por negocio. & Añade pruebas en CI y revisión por pares; aceptación del incremento en Review. \\
Control de calidad & Revisión de negocio en la columna de validación. & Parejas, pruebas y revisión técnica durante la implementación. \\
Mejora & Decisiones del patrocinador en el hito de control. & Retrospectivas con responsables y seguimiento, dentro del gobierno del Charter. \\
Review con cliente & Presentación de avance, riesgos y tablero. & Demo del incremento y registro de aceptación y observaciones. \\
\bottomrule
\end{tabularx}}

La simulación permitió contrastar el compromiso de 20 puntos con los 16 de S1. Las pruebas y revisiones
aportan evidencia del cierre técnico de 15 puntos. Las retrospectivas asignan responsables a las mejoras.
PMBOK conserva el control del presupuesto y de los cambios de alcance.

\section{Pronóstico para Sprint 3}
La propuesta prioriza US-06, modelo base de mora (13 puntos), con ventanas temporales, dataset persistido
y validación fuera de tiempo; el Charter fija AUC objetivo de 0,75. Churn (cinco puntos) requiere resolver
acceso y negociar alcance: ambos ítems completos suman 18 frente a la referencia de 15--16. El traslado
propuesto de US-07 a S4 exige revisar M4 con el patrocinador. En la Review se confirmará esta prioridad
y se completará \texttt{docs/SPRINT\_REVIEW\_ACCEPTANCE.md}. El guion prevé 900 segundos, con 150 de demo y
60 de preguntas. \mbox{README} vincula cada requisito con su solución y evidencia; la preparación actualizada
está en la PR \#9, pendiente de aprobación por otro integrante.

\end{document}
